\documentclass[11pt,a4paper]{article}

\usepackage[margin=2.5cm]{geometry}
\usepackage[T1]{fontenc}
\usepackage{lmodern}
\usepackage{microtype}
\usepackage{booktabs}
\usepackage{longtable}
\usepackage{xcolor}
\usepackage{listings}
\usepackage{amsmath}
\usepackage{enumitem}
\usepackage[hidelinks]{hyperref}

\definecolor{codebg}{RGB}{247,247,247}
\definecolor{codekey}{RGB}{20,80,160}
\definecolor{codestr}{RGB}{150,60,20}
\definecolor{codecom}{RGB}{110,110,110}

\lstdefinestyle{json}{
  basicstyle=\ttfamily\footnotesize,
  backgroundcolor=\color{codebg},
  breaklines=true,
  breakatwhitespace=false,
  columns=fullflexible,
  keepspaces=true,
  showstringspaces=false,
  frame=single,
  framerule=0pt,
  framesep=6pt,
  xleftmargin=6pt,
  stringstyle=\color{codestr},
  literate=
    *{:}{{{\color{codekey}:}}}{1}
     {\{}{{{\color{codekey}\{}}}{1}
     {\}}{{{\color{codekey}\}}}}{1},
}

\lstdefinestyle{plain}{
  basicstyle=\ttfamily\footnotesize,
  backgroundcolor=\color{codebg},
  breaklines=true,
  columns=fullflexible,
  keepspaces=true,
  showstringspaces=false,
  frame=single,
  framerule=0pt,
  framesep=6pt,
  xleftmargin=6pt,
}

\lstset{style=json}

\newcommand{\code}[1]{\texttt{\small #1}}
\newcommand{\file}[1]{\texttt{\small #1}}

% ---- GitHub links to files in the repository ---------------------------------
\newcommand{\ghbase}{https://github.com/ArashMassoudieh/OpenHydroQual/blob/master}
\newcommand{\ghtree}{https://github.com/ArashMassoudieh/OpenHydroQual/tree/master}
\newcommand{\exfile}[1]{\href{\ghbase/#1}{\file{#1}}}
\newcommand{\exdir}[1]{\href{\ghtree/#1}{\file{#1}}}

\title{\textbf{Composite Blocks in OpenHydroQual}\\[0.3em]
       \large A guide to building multi-block components}
\author{}
\date{}

\begin{document}
\maketitle
\vspace{-2.5em}

\begin{abstract}
\noindent
A \emph{composite} is a group of blocks and the links between them that the
user sees and manipulates as a single block. The computational engine is
unaware of composites: their members are ordinary \code{Block} and \code{Link}
objects, and every solver, parameter-estimation and output path treats them
individually. What a composite changes is the model file, which carries one
line instead of many, and the diagram, which draws one node instead of many.
This document explains how to declare a composite type in a plugin file, what
the engine does with that declaration, and the pitfalls that are easy to hit
and hard to diagnose.
\end{abstract}

\tableofcontents

\section{Concept}

A composite type is declared in an ordinary template (plugin) JSON file, next
to the block and link types, and is loaded the same way with
\code{addtemplate}. Dropping one on the diagram creates its member blocks and
the links between them in a single step; the members are then driven entirely
by the composite's own properties.

Three principles govern the design, and most of the practical advice in this
document follows from them.

\begin{description}[style=nextline,leftmargin=1.2em]
\item[The engine never sees a composite.]
Members are real blocks and links in the system's normal containers. They are
assembled into the state-variable matrices exactly like anything else. A
composite therefore cannot change the physics, only how a model is written
down and drawn.

\item[Members are strictly encapsulated.]
A member's property values are fully determined by the composite's exposed
properties. Members are never edited individually: their property panels are
read-only, and any value the composite type does not expose is unreachable. If
a user needs a knob the type does not offer, they \emph{ungroup} the composite,
after which it is plain blocks and links.

\item[The composite is the unit of saving and of drawing, not of computing.]
The model file carries \code{create composite;} with the group-level property
values; the members are regenerated from the type when the file is read. In the
diagram a composite is one node.
\end{description}

\section{Anatomy of a declaration}

A composite type is a top-level entry in a template JSON file with
\code{"type": "composite"}. Beyond the usual header fields it recognises three
reserved keys --- \code{members}, \code{internal\_links} and
\code{external\_links} --- and every other entry is an ordinary quantity
declaration that may carry an \code{applyto} mapping.

\begin{lstlisting}
"Clarifier": {
  "type": "composite",
  "typecategory": "Blocks",
  "description": "Clarifier: upper and lower settling compartments",
  "icon": { "filename": "clarifier.png" },
  "helptext": "Influent and effluent attach to the upper compartment; ...",

  "members": {
    "Upper": { "type": "Settling compartment",
               "dx": "0", "dy": "0",   "_width": "200", "_height": "200" },
    "Lower": { "type": "Settling compartment",
               "dx": "0", "dy": "250", "_width": "200", "_height": "200" }
  },

  "internal_links": {
    "interface": { "type": "Sedimentation tank element interface",
                   "from": "Upper", "to": "Lower" }
  },

  "external_links": {
    "*":              { "from": "Upper", "to": "Upper" },
    "Effluent flow":  { "from": "Upper", "to": "Upper" },
    "Sludge flow":    { "from": "Lower", "to": "Lower" }
  },

  "surface_area": {
    "type": "value",
    "description": "Cross-sectional area of the compartment interface",
    "ask_user": "true", "delegate": "UnitBox", "unit": "m^2;ft^2",
    "default": "100", "estimate": "true",
    "criteria": "surface_area>0",
    "warningmessage": "The surface area must be positive!",
    "applyto": { "interface#area": "surface_area" }
  }
}
\end{lstlisting}

\subsection{Header fields}

\begin{longtable}{@{}p{0.22\textwidth}p{0.72\textwidth}@{}}
\toprule
\textbf{Key} & \textbf{Meaning}\\
\midrule
\endhead
\code{type} & Must be \code{"composite"}.\\
\code{typecategory} & Toolbar and object-browser grouping. \code{"Blocks"} puts
the composite alongside ordinary blocks, which is usually what you want since
the user treats it as one.\\
\code{description} & Shown as the toolbar tooltip and the menu entry.\\
\code{icon} & \code{\{"filename": "name.png"\}}, resolved against
\exdir{resources/Icons}. See section~\ref{sec:register}.\\
\code{helptext} & Long-form help for the component.\\
\bottomrule
\end{longtable}

\subsection{Members}

Each entry declares one member block. The key is the \emph{member role} and
becomes part of the instance name.

\begin{longtable}{@{}p{0.16\textwidth}p{0.78\textwidth}@{}}
\toprule
\textbf{Field} & \textbf{Meaning}\\
\midrule
\endhead
\code{type} & The block type to instantiate. It must exist in the metamodel
when the composite is used, so make sure the template that defines it is loaded
--- either the same file, or one the user also adds.\\
\code{dx}, \code{dy} & Position \emph{offset} from the composite's own
\code{x}/\code{y}. Members are not drawn, but these offsets decide the layout
after an ungroup, so give them sensible values.\\
\code{\_width}, \code{\_height} & Member node size, used after ungrouping.\\
\bottomrule
\end{longtable}

\subsection{Internal links}

Each entry declares one link between two members. \code{from} and \code{to} are
member \emph{keys}, not instance names.

\begin{lstlisting}
"internal_links": {
  "percolation_1_2": { "type": "soil_to_soil_link",
                       "from": "Soil_1", "to": "Soil_2" }
}
\end{lstlisting}

\noindent
Internal links are part of the component: they are created with it, deleted
with it, and never appear in the diagram. Section~\ref{sec:internalreq} covers
an important obligation they carry.

\subsection{External links (ports)}\label{sec:ports}

When a user draws a link onto a composite node, the endpoint has to be resolved
to one of the members. Resolution is \emph{by link type}:

\begin{lstlisting}
"external_links": {
  "*":                        { "from": "Catchment",   "to": "Catchment" },
  "groundwater_link":         { "from": "Groundwater", "to": "Groundwater" },
  "groundwater_to_fixedhead": { "from": "Groundwater" }
}
\end{lstlisting}

\noindent
The key is a link type name; \code{"*"} is the fallback for any type not listed
explicitly. \code{from} applies when the composite is the link's source,
\code{to} when it is the destination. Omitting one direction means a link of
that type may not attach that way, and the user gets a clear message rather
than a silent mis-connection.

Once the port resolves, the ordinary link validation runs against the
\emph{member}, so a link that requires certain quantities on its endpoints is
still checked properly.

\paragraph{The one-type-one-port rule.}
A given link type can attach at only one place per direction. This is the main
constraint the port design imposes, and it shapes component design. In a
clarifier, influent and effluent attach to the upper compartment while return
and waste sludge leave the lower one --- if all four were the same generic
\code{Fixed flow} type, they could not be told apart. The remedy is to give the
distinct roles distinct link types (\code{Effluent flow} and
\code{Sludge flow}), which is also self-documenting in the diagram. Declaring
such a type is cheap: it is usually a copy of the generic flow link with a
different name, colour and icon.

\subsection{Exposed properties and mappings}\label{sec:mappings}

Every remaining entry is an ordinary quantity declaration --- the same fields
you would write for a block property (\code{type}, \code{description},
\code{ask\_user}, \code{delegate}, \code{unit}, \code{default},
\code{estimate}, \code{criteria}, \code{warningmessage}, \code{helptext}) ---
plus an optional \code{applyto} object that says which member quantities it
drives.

\begin{lstlisting}
"depth_to_groundwater": {
  "type": "value",
  "description": "Depth from the ground surface to the water table",
  "ask_user": "true", "delegate": "UnitBox", "unit": "m;ft",
  "default": "3.5", "estimate": "true",
  "applyto": {
    "Soil_1#depth": "depth_to_groundwater/7",
    "Soil_2#depth": "2*depth_to_groundwater/7",
    "Soil_3#depth": "4*depth_to_groundwater/7"
  }
}
\end{lstlisting}

Each \code{applyto} entry is \code{"<member>\#<quantity>": "<mapping>"}.

\begin{itemize}[leftmargin=1.4em]
\item \textbf{The target} names a member block \emph{or an internal link} ---
both are addressable, so an internal link's own properties can be driven from
the group level.
\item \textbf{The separator} between object and quantity is \code{\#}. It is
not \code{.} (reserved by the expression parser for the \code{.s}/\code{.e}
endpoint suffixes) and not \code{:} (already the constituent scoping separator
inside quantity names).
\item \textbf{One property may drive many targets}, and many properties may
drive targets on the same member. All mappings are re-applied together whenever
propagation runs, so it does not matter which property hosts which mapping ---
group them wherever they read most clearly.
\end{itemize}

\subsubsection{How a mapping is evaluated}\label{sec:eval}

This is the single most important rule in the document, and getting it wrong
produces confusing errors.

\begin{description}[style=nextline,leftmargin=1.2em]
\item[Numeric targets --- types \code{value}, \code{balance}, \code{constant}.]
The mapping is an arithmetic \emph{expression}, evaluated against the
composite's own properties. Anything the expression engine understands is
allowed.
\item[Every other target --- \code{source}, \code{timeseries}, \code{string},
\code{boolean}, \code{expression}.]
The mapping must be the \textbf{bare name of a composite property}; its value
is passed through unchanged. Arithmetic is not available here.
\end{description}

\noindent
Note that \code{expression}-typed quantities fall in the second group even
though they hold numbers. Several common initial conditions are declared this
way --- \code{Soil.theta} and \code{Groundwater cell.moisture\_content} are
both \code{expression} --- so they must be driven by a bare property name:

\begin{lstlisting}
"initial_moisture_content": {
  "type": "value", "ask_user": "true", "delegate": "ValueBox",
  "default": "0.2",
  "applyto": {
    "Soil_1#theta": "initial_moisture_content",
    "Soil_2#theta": "initial_moisture_content",
    "Soil_3#theta": "initial_moisture_content"
  }
}
\end{lstlisting}

\noindent
If you need arithmetic on such a target --- say a saturation fraction times a
porosity --- expose the derived value itself instead, and document the
relationship in the help text.

\paragraph{Expressions see only composite properties.}
A mapping expression is evaluated against the composite, so every symbol in it
must be another composite-level property. It cannot reach into a member. If a
formula needs a member's parameter, expose that parameter at the group level
too and drive the member from it.

\subsection{Naming}

\begin{longtable}{@{}p{0.34\textwidth}p{0.60\textwidth}@{}}
\toprule
\textbf{Form} & \textbf{Example}\\
\midrule
\endhead
Member instance name & \code{Clarifier\_1\_\_Upper} \quad
(\code{<instance>\_\_<member key>})\\
Internal link instance name & \code{Clarifier\_1\_\_interface}\\
Mapping target & \code{Upper\#Storage}\\
Output column & \code{Clarifier\_1\_\_Upper\_Storage} \quad
(\code{<block>\_<quantity>}, unchanged from ordinary blocks)\\
\bottomrule
\end{longtable}

\noindent
Because \code{\_\_} is the instance separator, avoid it in instance names you
type by hand. Member keys become user-visible in the object browser and the
results menu, so choose readable ones: \code{Upper}, \code{Lower},
\code{Soil\_1}, \code{interface}.

\section{What happens at run time}

\subsection{Instantiation}

When a composite is created --- from the toolbar, from a file, or by pasting
--- the following happens in order:

\begin{enumerate}[leftmargin=1.6em]
\item member blocks are created, named \code{<instance>\_\_<key>};
\item internal links are created and connected between them;
\item geometry is applied: each member's \code{x}/\code{y} is the composite's
position plus the declared offsets, and the members' geometry quantities are
hidden from the property panel;
\item constituent- and reaction-derived properties are exposed on the composite
(section~\ref{sec:constituents});
\item the group-level values from the file are assigned;
\item \textbf{propagation} runs: every \code{applyto} mapping is evaluated and
written to its target.
\end{enumerate}

\noindent
Members are created in alphabetical order of their keys, so if creation order
ever matters to you, name the keys accordingly.

\subsection{When propagation runs}

Mappings are re-applied at three moments, and it is worth knowing all three:

\begin{enumerate}[leftmargin=1.6em]
\item at instantiation, as above;
\item when the user edits a group-level property in the diagram or the property
panel;
\item during \code{ApplyParameters()}, so that a calibrated value reaches the
members.
\end{enumerate}

\noindent
The third is what makes parameter estimation work through a composite, and it
has a corollary: \textbf{parameters bind at the composite level only}. A
parameter assigned directly to a member quantity would be overwritten by the
next propagation, so the property panel disables parameter assignment for
members.

\subsection{Constituents and reactions}\label{sec:constituents}

When a constituent, reaction or reaction parameter is added, the engine copies
a set of derived quantities onto every block --- \code{X\_bh:concentration} and
friends. Their names are only known at run time, so a composite type cannot
declare mappings for them.

They are therefore exposed automatically: for each derived quantity found on
one or more members, a group-level property of the same name is created that
fans out to \emph{every} member carrying it. All members consequently share one
value for such quantities. If your type declares a property with the same name
explicitly, your declaration wins and no automatic mapping is generated.

\section{Persistence}

A composite is written as a single command carrying only its group-level
properties:

\begin{lstlisting}[style=plain]
create composite;type=Clarifier,name=Clarifier_1,x=1273,y=994,
    surface_area=100[m~^2],underflow=95[m~^3/day],upper_storage=100[m~^3], ...
\end{lstlisting}

\noindent
No members or internal links appear in the file --- they are regenerated from
the type when it is read. External links attached to a composite are written
against the \emph{composite} name and re-resolved to a member by link type on
load:

\begin{lstlisting}[style=plain]
create link;from=Clarifier_1,to=Effluent,type=Effluent flow,flow=95,name=...
\end{lstlisting}

\noindent
A consequence worth stating plainly: a model file no longer fully describes its
own structure. Editing a composite type in the template changes every saved
model that uses it. Treat a published composite type as an interface.

\paragraph{Expanded output.}
Both serialisers can also write a model in \emph{expanded} form, in which
composites are omitted and their members and internal links are written as
ordinary blocks and links --- the same output an ungrouped model produces. This
is available as an API flag (\code{expand\_composites}) for headless use and
export; the GUI always saves in composite form. Note that a parameter bound to
a group-level property has nowhere to live in an expanded file, so calibration
setup is not preserved there.

\section{Behaviour in the interface}

\begin{longtable}{@{}p{0.26\textwidth}p{0.68\textwidth}@{}}
\toprule
\textbf{Action} & \textbf{Result}\\
\midrule
\endhead
Toolbar / Blocks menu & Composite types appear after the block types, with
their icon and description.\\
Diagram & One node, drawn with the composite's icon. Members and internal links
are never drawn.\\
Moving / resizing & Members follow, keeping their declared offsets.\\
Property panel & Selecting the node shows the group-level properties. Member
geometry (\code{x}, \code{y}, \code{\_width}, \code{\_height}) is hidden, since
members have no node of their own.\\
Object browser & The composite appears with its members and internal links
nested beneath it, read-only, so member values stay inspectable.\\
Results menu & Grouped one submenu per member and internal link, listing that
member's output items.\\
Ungroup & Right-click the node. Members and links become independent objects
and the model saves flat afterwards. One-way: there is no regroup.\\
Delete & Removes the composite, its members, its internal links and any
external link attached to them.\\
Copy / paste & Copies the composite and its internal links. External links are
not copied.\\
\bottomrule
\end{longtable}

\section{Registering a plugin}\label{sec:register}

Put the composite in any template JSON in \file{resources/}. To make the file
selectable from the plugin chooser, add a line to
\exfile{resources/default\_templates.list}:

\begin{lstlisting}[style=plain]
Hydrologic Response Unit, hydrologic_response_unit.json, hydrologic_response_unit.png
\end{lstlisting}

\noindent
The three fields are the display name, the JSON file and the icon. Icons live
in \exdir{resources/Icons} and follow the existing convention: 256$\times$256
PNG with a transparent background, cropped to the artwork and centred. The same
image may be used for the plugin entry and for the composite type.

\section{Pitfalls}

Every item below was encountered while building the two shipped composites.
They share a characteristic: the model still loads and runs, so nothing draws
attention to the mistake.

\subsection{Drive the real initial-condition handle}\label{sec:initial}

Many storage quantities are derived, not set. \code{Soil.Storage} is declared
\code{ask\_user: false} with
\code{initial\_value\_expression: "area*depth*theta"}, and
\code{Groundwater cell.Storage} likewise derives from
\code{area*depth*moisture\_content}. Mapping such a \code{Storage} directly
appears to work --- the value is written --- but it is recomputed from its
expression during initialisation and your value is discarded.

\begin{quote}
\textbf{Symptom:} changing the exposed property has no effect whatsoever; two
runs with different initial conditions agree to many decimal places.
\end{quote}

\noindent
Map the quantity the expression reads --- \code{theta},
\code{moisture\_content} --- not the storage it produces. Both happen to be
\code{expression}-typed, so by section~\ref{sec:eval} their mappings must be
bare property names.

\subsection{Satisfy the internal links' required properties}\label{sec:internalreq}

An internal link belongs to the composite, so nobody else will configure it. If
its type declares a property with \code{ask\_user: true} and no usable default,
the composite \emph{must} map it. \code{soil2groundwater\_link.area} is exactly
this case: no default, so it starts at zero.

\begin{quote}
\textbf{Symptom:} a flux that should exist is identically zero --- in that
example, no water ever reaches the groundwater --- while the model solves
happily.
\end{quote}

\noindent
Before shipping a type, list every \code{ask\_user} property of every internal
link type and confirm each is either mapped or has a sound default.

\subsection{Anything unexposed is unreachable}

This is encapsulation working as designed, but it puts real weight on the
author. A property the type does not expose cannot be edited, and cannot even
be saved --- it reverts to the type's default on reload. Expose everything a
user might reasonably want to change, and remember that this includes initial
conditions. The escape valve is ungrouping.

\subsection{Watch out for confined-style head formulations}

Some block types compute head in a way that punishes a plausible-looking
initial condition. \code{Groundwater cell} uses
\[
  \text{head} = \text{cell top} +
  \frac{\text{moisture content} - \text{porosity}}{\text{specific storage}},
\]
so at full saturation the head sits at the top of the cell, and a small deficit
becomes a large elastic underpressure: with a specific storage of $0.01$, a
moisture content $0.025$ below porosity puts the head $2.5$~m low. Choose
defaults that correspond to a sensible state, and say so in the help text.

\subsection{Port ambiguity}

Covered in section~\ref{sec:ports}: one link type attaches at one place per
direction. If two roles on your component need the same generic link type,
introduce a distinct type for one of them.

\section{A recipe}

\begin{enumerate}[leftmargin=1.6em]
\item Decide the members and the internal links, and check that the link types
you need exist and connect those block types.
\item For each internal link type, list its \code{ask\_user} properties and plan
a mapping or confirm a sound default (section~\ref{sec:internalreq}).
\item Decide the ports. Group the external link types by which member they
should attach to, and add a distinct link type if two roles collide
(section~\ref{sec:ports}).
\item Decide the exposed properties. Work outward from what a user would
actually type: dimensions, material properties, sources, initial conditions.
For each, find the member quantity that genuinely controls it
(section~\ref{sec:initial}).
\item Write the mappings, respecting the numeric / bare-name rule
(section~\ref{sec:eval}).
\item Give every property a description, a unit, a default, and where it helps
a \code{criteria} with a \code{warningmessage}.
\item Add the icon and register the plugin (section~\ref{sec:register}).
\item Test (section~\ref{sec:testing}).
\end{enumerate}

\section{Testing a new composite}\label{sec:testing}

Build a small model that uses the component and check, in order:

\begin{enumerate}[leftmargin=1.6em]
\item \textbf{It instantiates.} The expected number of blocks and links appear,
with the expected names.
\item \textbf{Derived values are right.} Read a few member quantities back and
compare against what the mappings should have produced --- especially anything
computed from several group properties.
\item \textbf{Every knob bites.} Change each exposed property and confirm a
member value actually changes. This is what catches the
\hyperref[sec:initial]{initial-condition trap}, and nothing else does.
\item \textbf{Ports resolve.} Draw one external link of each type and confirm it
lands on the intended member.
\item \textbf{It round-trips.} Save, reload, and confirm the member values and
port resolutions survive.
\item \textbf{It verifies and solves.} \code{VerifyAllQuantities} should report
nothing, and the run should complete. Check the result is physically sensible,
not merely finite.
\item \textbf{Ungroup works.} The members should survive as plain blocks with
their geometry restored.
\end{enumerate}

\section{Worked example: a two-compartment clarifier}

The full type is in \exfile{resources/wastewater.json}. Two
\code{Settling compartment} members joined by a
\code{Sedimentation tank element interface}; influent and effluent attach to
the upper compartment, return and waste sludge leave the lower one. Because all
four would otherwise be \code{Fixed flow}, two dedicated link types are
declared alongside it, \code{Effluent flow} and \code{Sludge flow}, which give
the ports something to key on and make the diagram readable. Group properties
cover the interface area, the underflow, the two compartment elevations, the
initial volumes, and an external inflow series.

A model using it is in \exfile{Examples/ASM\_Clarifier\_Composite.ohq}: an
activated-sludge plant of four reactors whose settler is a single node.

\section{Worked example: a hydrologic response unit}

The full type is in \exfile{resources/hydrologic\_response\_unit.json}. Five
members --- a \code{Catchment}, three \code{Soil} layers and a
\code{Groundwater cell} --- with infiltration, two percolation links and a
recharge link between them. It illustrates several ideas at once:

\begin{itemize}[leftmargin=1.4em]
\item \textbf{Derived geometry.} The user supplies a surface elevation, a depth
to groundwater and a groundwater thickness; the layer thicknesses (in a
$1\!:\!2\!:\!4$ ratio, finest at the surface where moisture gradients are
steepest) and all four bottom elevations are computed:
\begin{lstlisting}
"Soil_1#bottom_elevation": "surface_elevation-depth_to_groundwater/7",
"Soil_2#bottom_elevation": "surface_elevation-3*depth_to_groundwater/7",
"Soil_3#bottom_elevation": "surface_elevation-depth_to_groundwater"
\end{lstlisting}
\item \textbf{One property, many members.} A single \code{area} drives the
catchment, all three soil layers, the groundwater cell \emph{and} the recharge
link's interface area.
\item \textbf{Selective application.} Evapotranspiration maps to the top soil
layer only, while precipitation goes to the catchment.
\item \textbf{Shared material properties.} The three layers discretise one soil
profile rather than representing different materials, so they share one
\code{K\_sat}, one \code{alpha\_vG} and so on. Splitting these per layer is a
mechanical edit if heterogeneous layers are wanted.
\end{itemize}

\noindent
A two-unit hillslope model using it is in
\exfile{Examples/Hydrologic\_Response\_Unit/HRU\_hillslope.ohq}.

\section{Limitations}

\begin{itemize}[leftmargin=1.4em]
\item \textbf{The member list is fixed.} A composite type declares a definite
set of members; the count cannot be driven by a property. Variable-length
structures --- an $n$-layer column, a channel reach of $n$ segments, an
$n_x\times n_y$ aquifer grid --- need one type per size, or the existing grid
generator.
\item \textbf{No nesting.} A composite cannot contain another composite.
\item \textbf{No per-member unique data.} Values that are neither shared nor a
formula over group properties --- a measured conductivity for each individual
layer, say --- cannot be expressed. Such models are built by ungrouping.
\item \textbf{Ungroup is one-way.} There is no regroup.
\end{itemize}

\end{document}
